% Part: computability
% Chapter: computability-theory
% Section: introduction

\documentclass[../../../include/open-logic-section]{subfiles}

\begin{document}

\olfileid{cmp}{thy}{int}
\olsection{परिचय}

\emph{संगणनक्षमता सिद्धांत} म्हणून ओळखली जाणारी तर्कशास्त्राची शाखा
फलने व संच यांची संगणनक्षमता किंवा सापेक्ष संगणनक्षमता यांच्याशी संबंधित
प्रश्नांचा अभ्यास करते. या विषयाला पूर्वी \emph{पुनरावर्तन सिद्धांत}
म्हणत असत, यावरून क्लिनींचा प्रभाव दिसतो; आज ही दोन्ही नावे
सर्वसाधारणपणे वापरली जातात.

एखाद्या संगणनाच्या प्रतिमानात फलन~$f\colon \Nat \pto \Nat$ चे संगणन
करता येत असेल, तर त्याला \emph{आंशिक संगणनक्षम} म्हणू. $f$ सर्वत्र
परिभाषित असेल, तर $f$ ला फक्त \emph{संगणनक्षम} म्हणू. ज्याचे
जातिबोधक फलन~$\Char{R}$ संगणनक्षम आहे, असा संबंध~$R$ देखील
संगणनक्षम म्हणतात. $f$ व~$g$ ही आंशिक फलने असतील, तर
$f(x) \fdefined$ याचा अर्थ $f$ हे~$x$ येथे परिभाषित आहे, म्हणजे
$x$ हे~$f$ च्या परिभाषाक्षेत्रात आहे; $f(x) \fundefined$ याचा
अर्थ उलट, म्हणजे $f$ हे~$x$ येथे परिभाषित नाही. $f(x) \simeq g(x)$
याचा अर्थ $f(x)$ आणि $g(x)$ ही दोन्ही अपरिभाषित आहेत, किंवा दोन्ही
परिभाषित असून समान आहेत.

संगणनाच्या एखाद्या विशिष्ट प्रतिमानाचा उल्लेख न करताही संगणनक्षमता
सिद्धांताचा अभ्यास करता येतो. त्यासाठी एक सार्वत्रिक आंशिक संगणनक्षम
फलन~$\fn{Un}(k, x)$ आहे, हे दाखवतात. त्याच्या आधारे आंशिक संगणनक्षम
फलनांचे प्रगणन करता येते. $k$-वे एकस्थानी आंशिक संगणनक्षम फलन दर्शवण्यासाठी
आपण~$\cfind{k}$ हे चिन्हांकन वापरू; त्याची व्याख्या
$\cfind{k}(x) \simeq \fn{Un}(k, x)$ अशी आहे. (क्लिनी यांनी यासाठी
$\{ k \}$ हे चिन्हांकन वापरले होते; पण अलीकडे ते तितकेसे वापरले
जात नाही.) यापेक्षा थोड्या अधिक सामान्य स्वरूपात, कोणत्याही
स्थानसंख्येच्या आंशिक संगणनक्षम फलनांचे एकसमान प्रगणन करता येते.
$k$-वे $n$-स्थानी आंशिक पुनरावर्ती फलन दर्शवण्यासाठी आपण
$\cfind{k}[n]$ वापरू.

$f(\vec x, y)$ हे सर्वत्र परिभाषित किंवा आंशिक फलन असेल, तर
$\umin{y}{f (\vec x, y)}$ हे~$\vec x$ चे असे फलन आहे जे
$f(\vec x, y) = 0$ होईल असा लघुतम~$y$ परत करते, मात्र
$f(\vec x, 0)$, \dots, $f(\vec x, y-1)$ ही सर्व मूल्ये परिभाषित
असली पाहिजेत. असा~$y$ नसल्यास $\umin{y}{f (\vec x, y)}$
अपरिभाषित असते. $R(\vec x, y)$ हा संबंध असेल, तर
$\umin{y}{R(\vec x, y)}$ ची व्याख्या $R(\vec x, y)$ सत्य
असेल असा लघुतम~$y$ अशी केली जाते; दुसऱ्या शब्दांत,
$1 \tsub \Char{R}(\vec x, y) = 0$ होईल असा लघुतम~$y$.

एखादे फलन संगणनक्षम आहे हे दाखवण्यासाठी पुढील दोन मार्ग आहेत:
\begin{enumerate}
\item काटेकोरपणे: ट्यूरिंग यंत्र किंवा आंशिक पुनरावर्ती फलन स्पष्टपणे
  वर्णन करणे आणि ते इच्छित फलनाचे संगणन करते हे दाखवणे;
\item अनौपचारिकपणे: त्याचे संगणन करणाऱ्या अल्गोरिदमचे वर्णन करणे आणि
  चर्च यांच्या प्रबंधाचा आधार घेणे.
\end{enumerate}
या दोन मार्गांमध्ये नेमकी सीमारेषा नाही. अल्गोरिदमचे तपशीलवार
वर्णन पुरेशी माहिती देणारे असावे, जेणेकरून तत्त्वतः योग्य ट्यूरिंग
यंत्र किंवा आंशिक पुनरावर्ती व्याख्यांची मालिका कशी तयार करायची हे
तुलनेने स्पष्ट होईल. पूर्णपणे काटेकोर व्याख्या फारशा माहितीपूर्ण
असण्याची शक्यता नाही. म्हणून आपण या दोन मार्गांमध्ये योग्य मध्यममार्ग
शोधू; थोडक्यात, अचूक व्याख्या मिळवता येतील याची खात्री देणारे
सहज समजणारे पण काटेकोर पुरावे शोधू.

\end{document}
